% Slide 2 – the problem as a cost for three parties, and the market we start with
% (criteria: relevance 25 %, business model 20 %; RULES: Idea).
\begin{frame}{Who pays today for missing accessibility information}
  \begin{columns}[T,onlytextwidth]
    \begin{column}{0.56\textwidth}
      \begin{block}{Three parties, one data gap}
        \footnotesize
        \begin{itemize}
          \item \textbf{Guest} (wheelchair, pram, senior): does not know whether they can get in, so they do not book
                or make a wasted trip. An “accessible” label without a source or date says nothing.
          \item \textbf{Venue:} loses customers it never sees and answers the same questions about steps
                and toilets by e-mail and phone.
          \item \textbf{City:} has a statutory duty to ensure accessibility, but no current data
                and no capacity to maintain a database by hand.
        \end{itemize}
      \end{block}
      \begin{exampleblock}{Scope of the prototype}
        \footnotesize
        Wheelchair users and parents with prams. The same data serves 8~needs profiles, from seniors to
        medical tourists, and the interface works in Polish, English, German and Ukrainian.
      \end{exampleblock}
    \end{column}
    \begin{column}{0.42\textwidth}
      \textbf{\small The market we start with}\\
      \muted{Kraków, the Accessly catalogue built from OpenStreetMap}
      \vspace{1mm}

      \kpi{502}{accommodation}\hfill\kpi{2,622}{food \& drink}\hfill\kpi{1,117}{health}
      \vspace{2mm}

      \textbf{\small Demand}\textsuperscript{*}
      \footnotesize
      \begin{itemize}
        \item about 14~million visitors to Kraków a year
        \item 5.4~million people with disabilities in Poland; one resident in five is 65 or older
        \item growing medical tourism: patients from abroad ask about accessibility before booking
      \end{itemize}
      \vspace{1mm}
      \muted{\textsuperscript{*} External figures to verify before the pitch: visitor surveys of the
        Małopolska Tourist Organisation, Statistics Poland (2021 census, population).}
    \end{column}
  \end{columns}

  \note[item]{Open with the story: Anna, a wheelchair user, wants a coffee in the centre. Last year's map says “accessible”; today there is a step at the door. She loses her evening, the café loses a customer, and the city does not know the step exists.}
  \note[item]{The figures on the right come from our catalogue (420 culture venues too). Starred figures are public data that must be confirmed, with a source ready if the jury asks.}
  \note[item]{Primary group from the brief: wheelchairs and prams. The other profiles reuse the same attributes, so they widen the market without extra data.}
\end{frame}
